%% ma: L12-B07-0020
%% lop: 12
%% chuong: 2
%% bai: 7
%% dang: 12.7.9
%% loai: TN4
%% muc: NB
%% dap_an: "B"
%% nguon: "CĐ2.2.pdf giáo viên gửi 27/10/2026 (M07, Lê Minh Khôi – Nguyễn Anh Khôi, Chương II Chủ đề 2 Tọa độ của vectơ trong không gian), Câu 17"
%% kien_thuc: "Bài 7 (hình chiếu của điểm lên trục tọa độ), Bài 8 (trung điểm)"
%% trang_thai: chinh-thuc
%% ghi_chu: "Giáo viên đã duyệt 27/10/2026"
\begin{ex}
Trong không gian $Oxyz$, cho điểm $T(11;3;7)$. Gọi $T_1$ là điểm đối xứng của $T$ qua trục tung. Tọa độ của điểm $T_1$ là
\choice{$T_1(0;3;0)$}{\True $T_1(-11;3;-7)$}{$T_1(11;-3;7)$}{$T_1(0;-3;0)$}
\loigiai{Hình chiếu của $T$ lên $Oy$ là $H(0;3;0)$, là trung điểm $TT_1$ nên $T_1(-11;3;-7)$.}
\end{ex}
